%قبل از شروع مسئله ترجیحاً چیزی ننویسید. اما اگر خواستید بنویسید، در قالب لیستی (فقرات و شمارش) بنویسید.
\شروع‌مسئله[30] \متن‌سیاه{سؤال‌های جواب کوتاه.} در سؤال‌های این قسمت معمولاً یک جواب کوتاه و/یا یک توضیح کوتاه حدوداً بین یک تا دو خط مدنظر است؛ مگر این که صورت آن قسمت به صراحت خواستۀ دیگری داشته باشد.
% صورت اصلی مسئله را در این‌جا بنویسید.
\شروع{زیرمسائل}
\زیرمسئله[6]
این جمله درست است یا غلط؟
وقتی $n$ عنصر با کلیدهای اعداد صحیح متفاوت را در یک جدول درهم‌سازی ذخیره می‌کنیم که تابع درهم‌سازی آن به صورت تصادفی از یک خانواده توابع درهم‌سازی سراسری (\متن‌لاتین{universal hashing}) انتخاب شده باشد، برخوردی در جدول درهم‌سازی رخ نمی‌دهد.

\پاسخ{
غلط است. تمام توابع درهم‌سازی ممکن است کلیدهای متفاوت را به مقادیر یکسان نگاشت کند.

}

\زیرمسئله[6]
این جمله درست است یا غلط؟
در یک جدول درهم‌سازی که از درهم‌سازی کوکو استفاده می‌کند، زمان اجرای جست‌وجو در بدترین حالت برابر است با $O(1)$.

\پاسخ{
درست است زیرا جست‌وجو به حداکثر دو بار درهم‌سازی نیاز دارد.
	
}

\زیرمسئله[6]
چرا در تقسیم اعداد بزرگ با روش نیوتون، دقت محاسبات را رفته رفته افزایش می‌دهیم (به جای این‌که از ابتدا با بیشترین دقت محاسبات را انجام دهیم)؟

\پاسخ{
	زیرا در این صورت است که زمان تقسیم با ضرب به صورت مجانبی برابر می‌شود (اگر از ابتدا با بیشترین دقت محاسبات را انجام دهیم، زمان تقسیم $\Theta(\log n)$ برابر زمان ضرب خواهد بود).
	
}

\زیرمسئله[6]
برای تحلیل سرشکن آرایۀ پویا (که فقط عمل درج دارد)، یک تابع پتانسیل مناسب پیشنهاد دهید.

\پاسخ{
	اگر طول آرایه برابر با $N$ و تعداد عناصر موجود در آن $n$ باشد، یک تابع مناسب
	$(2n-N)c$
	می‌باشد که $c$ یک عدد ثابت مناسب است. در صورتی که ضرب در $c$ لحاظ نشود، نیمی از نمره داده می‌شود.
	
}

\زیرمسئله[6]
درخت دکارتی مربوط به دنباله اعداد زیر را رسم کنید.
\begin{equation*}
6, 9, 2, 4, 7, 8, 5, 8, 3, 7
\end{equation*}

\پاسخ{
	https://cp-algorithms.com/graph/rmq_linear.html
}

%\زیرمسئله[3]
%
%
%\پاسخ{
%
%
%	
%}
%
%\زیرمسئله[3]
%
%
%
%\پاسخ{
%
%
%	
%}

%%\begin{figure}[h!]%[5]{L}{0.4\textwidth}
%	%	\vspace{-2cm}
%	\begin{center}
%		\begin{tikzpicture}[-,>=stealth',level/.style={sibling distance = 5cm/#1,
%			level distance = 1cm}] 
%		\node [arn_s] {5}
%		child{ node [arn_s] {1} 
%			child{ node [arn_x] {}}
%			child{ node [arn_x] {}}
%		}
%		child{ node [arn_s] {15}
%			child{ node [arn_s] {10} 
%				child{ node [arn_s] {7}
%					child{ node [arn_x] {}}
%					child{ node [arn_x] {}}
%				}
%				child{ node [arn_s] {13}
%					child{ node [arn_x] {}}
%					child{ node [arn_x] {}}
%				}
%			}
%			child{ node [arn_s] {20}
%				child{ node [arn_x] {}}
%				child{ node [arn_x] {}}
%			}
%		}
%		; 
%		\end{tikzpicture}
%	\end{center}
%%	\caption{درخت قرمز-سیاه}
%%	\label{شکل:۷و۸}
%%\end{figure}



\پایان{زیرمسائل}



%اگر مسئله چند قسمت داشت به این صورت عمل کنید:
%\شروع{زیرمسائل}
%\زیرمسئله[0] % مسئلۀ ۱آ
%\زیرمسئله[0] % مسئلۀ ۱ب
%\زیرمسئله[0] % مسئلۀ ۱پ
%\زیرمسئله[0] % مسئلۀ ۱ت
%\پایان{زیرمسائل}





%%%%%%%%%%%%%%%%%%%%%%%%%%%%%%%%%%%%%%%%%%%%%%%%%%%%%
% برخی دستورات مفید برای تمرین‌ها. %
%%%%%%%%%%%%%%%%%%%%%%%%%%%%%%%%%%%%%%%%%%%%%%%%%%%%%

% فرمول‌های نمونه:
%$f(x) = \Theta(x)$
%$T(x, y) \leq \log(x) + 2^y + \binom{2n}{n}$
% {\tt code\_function}


% لیست موارد بدون شماره‌گذاری:
%\شروع{فقرات}
%    \فقره مورد اول 
%        \شروع{فقرات}
%            \فقره مورد اول داخلی 
%                \شروع{فقرات}
%                    \فقره اولین مورد خیلی داخلی 
%                    \فقره دومین مورد خیلی داخلی 
%                \پایان{فقرات}
%            \فقره مورد دوم داخلی 
%        \پایان{فقرات}
%    \فقره مورد دوم 
%\پایان{فقرات}

% لیست موارد شماره‌گذاری شده:
%\شروع{شمارش}
%    \فقره یک مورد 
%    \فقره یک مورد دیگر 
%    \فقره باز هم یک مورد دیگر
%\پایان{شمارش}

% برای نوشتن یک فرمول چند خطی که یک شماره بخوره. برای حالت‌های دیگر به سایت زیر مراجعه کنید:
% https://www.overleaf.com/learn/latex/Aligning_equations_with_amsmath
%\begin{align} 
%    \begin{split}
%        (x+y)^3 &= (x+y)^2(x+y) \\
%                &= (x^2+2xy+y^2)(x+y) \\
%                &= (x^3+2x^2y+xy^2) + (x^2y+2xy^2+y^3) \\
%                &= x^3+3x^2y+3xy^2+y^3
%    \end{split}                                 
%\end{align}

% نوشتن ماتریس و آرایه دوبعدی
%\begin{align}
%    A = 
%    \begin{bmatrix}
%        A_{11} & A_{21} \\
%        A_{21} & A_{22}
%    \end{bmatrix}
%\end{align}

% اضافه کردن تصویر با شرح. برای اضافه کردن تصویر بدون شرح، فقط دستور درج‌تصویر را استفاده کنید.
%\شروع{شکل}[h!]
%\تنظیم‌ازوسط
%  \درج‌تصویر[پهنا=0.5\پهنای‌متن]{img.jpg}
%\شرح{شرح تصویر}
%\برچسب{تصویر:۱}
%\پایان{شکل}
%\پایان{زیرمسائل}

%کد زیر را ببینید.
%%\lstinputlisting{source_filename.py}
%\begin{latin}
%\begin{lstlisting}
%# Import the modules
%import sys
%import random
%
%ans = True
%
%while ans:
%question = raw_input("Ask the magic 8 ball a question: (press enter to quit) ")
%
%answers = random.randint(1,8)
%
%if question == "":
%sys.exit()
%
%elif answers == 1:
%print "It is certain"
%
%elif answers == 2:
%print "Outlook good"
%
%elif answers == 3:
%print "You may rely on it"
%
%elif answers == 4:
%print "Ask again later"
%
%elif answers == 5:
%print "Concentrate and ask again"
%
%elif answers == 6:
%print "Reply hazy, try again"
%
%elif answers == 7:
%print "My reply is no"
%
%elif answers == 8:
%print "My sources say no"
%\end{lstlisting}
%\end{latin}
